\documentclass[11pt]{article}
\usepackage[T1]{fontenc}
\usepackage[utf8]{inputenc}
\usepackage{lmodern,amsmath,amssymb,amsthm}
\usepackage[margin=1in]{geometry}
\usepackage{hyperref}
\newtheorem{theorem}{Theorem}
\newtheorem{lemma}[theorem]{Lemma}
\newtheorem{proposition}[theorem]{Proposition}
\newtheorem{corollary}[theorem]{Corollary}
\theoremstyle{definition}
\newtheorem{definition}[theorem]{Definition}
\newcommand{\Cq}{\mathcal C_q}
\newcommand{\Iq}{\mathcal I_q}
\newcommand{\Prad}{P_{\rm rad}}
\newcommand{\ord}{\operatorname{ord}}
\title{Radial minorant envelopes on finite singular curve links\\
\large Partial progress on P12.3; the full atom remains open}
\author{Research draft for independent review}
\date{3 October 2026}
\begin{document}
\maketitle
\begin{abstract}
On an isolated curve singularity with disc normalization and a separated
conductor ideal, the known literal radial psh profiles are zero or strictly
negative on every branch. We rederive the classical part of this known
classification from Gubler--Rabinoff's primary source. We then compute the
radial obstacle supremum and show that mixed-zero obstacles have no maximal
radial minorant, although their supremum is continuous on the finite link.
No monotone minorant projection fixing this radial class can exist. An
explicit increasing family of logarithms on the node realizes the defect,
and positive diagonal changes of primary gauge preserve it. These are
consequences within a known restricted class, with no priority claim. They
do not solve any full active atom or construct products or finite energy.
\end{abstract}

\section{Scope, definitions, and external inputs}
Let $K$ be complete and nontrivially valued. Fix $\varpi\in K$ with
$0<|\varpi|<1$, and write $r=|\varpi|$, $a=-\log r>0$.
Let $A$ be a reduced strict $K$-affinoid algebra of pure dimension one,
with finite normalization
\[
 A\hookrightarrow B=\prod_{i=1}^qK\langle t_i\rangle,
 \qquad \nu:\coprod_{i=1}^q D_i\longrightarrow U=\mathcal M(A).
\]
Assume $\nu$ is finite and surjective, its fiber above a $K$-rational
point $o$ consists of the origins $p_i$, and it is an isomorphism away
from these origins. Assume explicitly that
\[
 \prod_{i=1}^qt_i^{e_i}K\langle t_i\rangle
 \subseteq \mathfrak c_{B/A}:=\{b\in B:bB\subseteq A\},
 \qquad e_i\geq1.
\]
These hypotheses define the partial range; they are not claimed for an
arbitrary excellent germ. In particular, there is no extension here to
arbitrary dimension, trivial valuation, or arbitrary normalization geometry.

For $c\in[-\infty,0]^q$, define a function off $o$ by
\[
 h_c^\circ(\nu(z))=
 \begin{cases}
  -(c_i/a)\log|t_i(z)|,&c_i> -\infty,\ z\in D_i\setminus\{p_i\},\\
  -\infty,&c_i=-\infty,\ z\in D_i\setminus\{p_i\}.
 \end{cases}
\]
A \emph{literal radial extension} fixes all these off-node values, not
just values at a finite set of points. At the radius-$r$ Gauss point
$\eta_i$ the value is $c_i$. We use the finite branch set
$L_q=\{1,\ldots,q\}$ with its discrete topology.

We use exactly classical psh in \cite[Definition 4.1]{GR}: an upper
semicontinuous function $u:X\to\mathbb R\cup\{-\infty\}$ is psh if
for every analytic extension $K'/K$, every $K'$-analytic curve $Y$, and
every morphism $f:Y\to X_{K'}$, its pullback through $f$ and the scalar
projection is subharmonic in Definition 3.4 of that source. The latter
includes the identically $-\infty$ convention. The essential external
statements, for any $K$-analytic space $X$, are:
\begin{enumerate}
\item For $f\in\Gamma(X,\mathcal O_X)$, $\log|f|$ is classically psh,
with $\log0=-\infty$; there is no nonzerodivisor condition
\cite[Proposition 4.2(2)]{GR}.
\item If $u_1,u_2$ are classically psh and $\lambda_1,\lambda_2\geq0$,
then $\lambda_1u_1+\lambda_2u_2$ and $\max(u_1,u_2)$ are classically
psh \cite[Proposition 4.2(3)]{GR}.
\item The infimum of a decreasing \emph{net} of classically psh functions
is classically psh \cite[Proposition 4.2(4)]{GR}.
\item For an analytic morphism $f:Y\to X$, pullback preserves classical
psh; the converse holds if $f$ is finite and surjective
\cite[Proposition 4.2(6)]{GR}.
\item A classically psh function on a $K$-affinoid space attains its
maximum at a Shilov point \cite[Theorem 6.4]{GR}.
\end{enumerate}
The source permits any complete absolute value in these statements.
All our pullbacks are through analytic morphisms, and all domains for
the maximum principle are affinoid closed discs. No boundaryless
maximum principle is used. Constants are psh by the definition. No CLD
comparison, current pullback, or wedge product is an external input here.

For completeness, a closed radius-$R$ disc has a singleton Shilov boundary
at its Gauss point: the triangle inequality bounds the value of
$\sum b_nt^n$ at every point by $\max_n|b_n|R^n$, and the Gauss seminorm
attains that bound for every series. This verifies the Shilov points used
below without an additional geometric theorem.

\section{A known classification rederived}
\begin{lemma}[Classical part of the known radial classification]
The profiles with a classical radial extension are exactly
\[
 \Cq^{\rm ext}=\{0\}\cup[-\infty,0)^q.
\]
The classically psh radial extension is unique, with value $0$ at $o$ for the zero profile and
$-\infty$ at $o$ otherwise. Its finite part is
$\Cq=\{0\}\cup(-\infty,0)^q$.
\end{lemma}
\begin{proof}
The zero and all-$-\infty$ profiles have their constant extensions.
For an all-negative profile let $I$ index its finite coordinates, and
put $\lambda_i=-c_i>0$ for $i\in I$. Choose positive rational numbers
$k_{n,i}/N_n\nearrow\lambda_i$ with a common denominator at each stage.
Multiplying the numerators and denominator by the same integer ensures
$k_{n,i}\geq e_i$. Put $b_{n,i}=t_i^{k_{n,i}}$ for $i\in I$ and
$b_{n,i}=0$ otherwise. The tuple belongs to the conductor ideal, hence
to $A$; denote it by $f_n$. Then
\[
 g_n=(aN_n)^{-1}\log|f_n|
\]
is classically psh by Proposition 4.2(2),(3), including when $f_n$ is a
zerodivisor. Its branch expressions are
$a^{-1}(k_{n,i}/N_n)\log|t_i|$ on retained branches and $-\infty$ on
the others; its value at $o$ is $-\infty$. Since $|t_i|\leq1$, this
sequence decreases to the desired extension. Proposition 4.2(4) applies.
This explicit rational sequence does not replace nets in that theorem.

Conversely, if $c_i=0$, the points $t_i=\varpi^m$ approach $p_i$.
Their images approach $o$, so upper semicontinuity of an extension $h$
forces $h(o)\geq0$. If some $c_j<0$, pull $h$ back to $D_j$ and
restrict it to the radius-$r$ disc. Theorem 6.4 gives
$h(o)\leq h(\eta_j)=c_j<0$, including $c_j=-\infty$, a contradiction.

For uniqueness when $c=0$, the same approach to $o$ gives $h(o)\geq0$
and the disc maximum principle gives $h(o)\leq0$. If some coordinate
$c_j$ is finite and negative, the radius-$r^m$ disc instead gives
$h(o)\leq mc_j$ for every $m$, forcing $-\infty$. If every coordinate
is $-\infty$, one disc maximum principle already gives that value.
All off-node values were prescribed.
\end{proof}

This lemma is the classical part of Book Theorem 70.1.9, already recorded
in the supplied catalogue \cite{Book}. We reproduce its argument to
check its external hypotheses. The book additionally identifies this
class with literal CLD-v2 radial extensions; that statement is not
reproved or used here and is not counted as new progress.

\section{The exact radial envelope and its defect}
\begin{definition}
For $\phi\in\mathbb R^q$, define the coordinatewise supremum
\[
 \Prad\phi:=\sup\{c\in\Cq:c\leq\phi\}.
\]
This definition does not presume that the output belongs to $\Cq$.
\end{definition}

\begin{theorem}\label{thm:envelope}
For every $\phi\in\mathbb R^q$,
\[
 \Prad\phi=\min(\phi,0)=:h.
\]
The supremum is a greatest radial minorant if and only if $h=0$ or
every $h_i<0$. If $h$ has both a zero and a negative coordinate, the
admissible set has no maximal element. Its supremum is continuous on
$L_q$ and equals its upper semicontinuous regularization, but is not in
$\Cq$. Allowing $\Cq^{\rm ext}$ does not change the supremum for a
finite obstacle.
\end{theorem}
\begin{proof}
Every admissible profile is bounded above by $h$. The profiles
$c^{(n)}=h-n^{-1}(1,\ldots,1)$ are strictly negative, are below
$\phi$, and increase uniformly to $h$. This proves the formula.
If $h\in\Cq$, it is itself the greatest minorant.

Suppose $h$ is mixed-zero. Zero is not admissible because $\phi$ has a
negative coordinate. Every admissible $c$ is therefore strictly negative.
Choose $0<\epsilon<\min_{h_i=0}(-c_i)$ and set
\[
 w_i=\begin{cases}h_i,&h_i<0,\\-\epsilon,&h_i=0.\end{cases}
\]
Then $w\in\Cq$, $c\leq w\leq h\leq\phi$, and $w$ strictly improves
every zero coordinate. Hence no $c$ is maximal. Finite functions on
the finite discrete set $L_q$ are continuous, so regularization cannot
alter $h$. The upper bound remains valid for extended admissible
profiles, and the same finite approximants give the lower bound.
\end{proof}

Thus the raw numerical formula agrees with the known ideal-orthant
envelope $\min(\phi,0)$, but the geometric radial codomain differs.
For $q=1$ no mixed-zero profile exists, consistently with the known
smooth curve and cusp tests. This is not a general smooth-germ theorem.

\begin{theorem}[No minorant projection]\label{thm:projection}
For $q\geq2$ no order-preserving map
$E:\Iq:=(-\infty,0]^q\to\Cq$ can satisfy both $E(\phi)\leq\phi$
for every $\phi$ and $E(c)=c$ for every $c\in\Cq$.
There is also no continuous retraction $\Iq\to\Cq$ fixing $\Cq$.
\end{theorem}
\begin{proof}
Choose a mixed-zero $h\in\Iq$. For $c^{(n)}=h-n^{-1}\mathbf1$,
monotonicity and the fixed-point condition give
$c^{(n)}=E(c^{(n)})\leq E(h)$. Taking coordinatewise suprema gives
$h\leq E(h)$; the minorant condition forces $E(h)=h\notin\Cq$.
For the second assertion, continuity and $c^{(n)}\to h$ force
$E(h)=\lim_n E(c^{(n)})=h$, again a contradiction. The first assertion
uses no continuity or limit-preservation axiom.
\end{proof}

The restrictions matter: there is an extensive monotone retraction
which fixes the zero and strictly negative profiles and sends every
mixed-zero profile to zero. If $c\leq d$ and $c$ is mixed-zero, then
$d$ has a zero coordinate too; this verifies monotonicity in the only
nontrivial case. Consequently we do not claim that every order map is
impossible. The obstruction concerns minorant envelopes and continuous
retractions.

\section{An explicit analytic realization on the node}
Take $A=K\langle x,y\rangle/(xy)$. Modulo $xy$, each series is uniquely
$\gamma+f(x)+g(y)$, with $f(0)=g(0)=0$. This embeds $A$ in
$B=K\langle x\rangle\times K\langle y\rangle$ as pairs with equal
constant terms; it proves reducedness and identifies the two pure
one-dimensional branches. The idempotents $(1,0),(0,1)$ generate $B$
as an $A$-module. The conductor is
$xK\langle x\rangle\times yK\langle y\rangle$: multiplication by
the two idempotents forces both constant terms to vanish, while such
pairs multiply $B$ into $A$. The analytic map is surjective with two
origins over $o$, and away from $o$ one branch coordinate is an analytic
unit locally, so the map is an isomorphism. This verifies all the analytic
hypotheses of the lemma with $e_x=e_y=1$.

For the local model $R=K[x,y]_{(x,y)}/(xy)$, a centered semivaluation with $0<v(\mathfrak m)<\infty$
has $v(x)+v(y)=\infty$. Both cannot be infinite because
$0<v(\mathfrak m)<\infty$. Exactly one is finite and positive. On the
corresponding one-variable quotient, every nonzero element is $t^d$
times a unit; units have value zero since they and their inverses have
nonnegative value. Thus every centered ray is a positive multiple of
$\ord_x$ or $\ord_y$, with their nonzero prime supports retained.
Normalizing by $\mathfrak m=(x,y)$ gives exactly these two points.
Evaluation of $x$ distinguishes $1$ from $\infty$, so their topology is
discrete. The same proof works for the formal node. The analytic and
link proofs use no excellence assertion; this is the known nodal range
of the catalogue, not a counterexample to the entire excellent-germ
programme.

Every element of $K^*$ is a unit in $R$, so the local semivaluations are
trivial on $K$. The analytic Gauss seminorms extend the nontrivial
absolute value of $K$. Only branch indices and the specified monomial
profiles are compared here; the two carriers are not identified.

\begin{proposition}\label{prop:node}
The classically psh functions
\[
 g_n=(an)^{-1}\log|x^n+y|
\]
increase on $U$ and have branch profiles $(-1,-1/n)$ at the two
radius-$r$ Gauss points. Their pointwise supremum is not upper
semicontinuous at $o$, and its usc regularization is not classically psh.
\end{proposition}
\begin{proof}
GR Proposition 4.2(2),(3) proves psh. On $D_x$ the function equals
$a^{-1}\log|x|$; on $D_y$ it equals $(an)^{-1}\log|y|$; at $o$ it is
$-\infty$. Since $\log|y|\leq0$, this increases. The Gauss values are
$(-1,-1/n)$, at uniform distance $1/n$ from $(-1,0)$.
The pointwise supremum $G$ is $a^{-1}\log|x|$ on $D_x\setminus\{o\}$,
zero on $D_y\setminus\{o\}$, and $-\infty$ at $o$. Its usc
regularization agrees away from $o$ and has value zero there: the
$y$-branch approaches $o$ with value zero, and all values are at most
zero. If this regularization were psh, its pullback to the closed
radius-$r$ $x$-disc would have central value zero and value $-1$ at
the unique Shilov point, violating Theorem 6.4.
\end{proof}

These are actual primary principal models: modulo $x^n+y$ and $xy$,
$x^{n+1}=y^2=0$, so $(x^n+y)$ is $\mathfrak m$-primary in $R$.
The functions are unbounded near $o$ on $U$, but finite and uniformly
bounded on the normalized link. This is not an obstruction to bounded
analytic Bedford--Taylor product continuity. GR Theorem 6.5 requires
the pointwise net limit already to be usc; $G$ is not usc and $G^*$ is
not that pointwise limit, so there is no contradiction to that theorem.

\section{Gauges and normalization compatibility}
For $\mathfrak m$-primary ideals $\mathfrak a,\mathfrak b$ on the node, let
$b_i^{\mathfrak a}=\ord_i(\mathfrak a)$,
$b_i^{\mathfrak b}=\ord_i(\mathfrak b)$. These numbers are finite and
positive because the ideals are proper and contain powers of
$\mathfrak m$. The section point is $\ord_i/b_i^{\mathfrak a}$.
For a degree-one homogeneous function $H$ on the rays,
\[
 c_i^{\mathfrak b}=H(\ord_i/b_i^{\mathfrak b})
 =\frac{b_i^{\mathfrak a}}{b_i^{\mathfrak b}}
    H(\ord_i/b_i^{\mathfrak a})
 =\frac{b_i^{\mathfrak a}}{b_i^{\mathfrak b}}c_i^{\mathfrak a}.
\]
This proof supplies the exact cocycle, since successive ratios telescope.
Positive diagonal scaling $T$ preserves both radial classes and
mixed-zero profiles and gives
$\Prad(T\phi)=T\Prad(\phi)$ by Theorem~\ref{thm:envelope}.
For $\mathfrak a=(x,y)$ and $\mathfrak b=(x^2,y^3)$,
$(-1,0)$ becomes $(-1/2,0)$ and still fails the radial condition.
These are changes of section, not identifications of all normalizers.
For the same literal analytic function its coefficient at the fixed
radius-$r$ Gauss points is
$d_i=H(\ord_i)=b_i^{\mathfrak a}c_i^{\mathfrak a}$, and the branch
expression is $-(d_i/a)\log|t_i|$. The identity
$b_i^{\mathfrak a}c_i^{\mathfrak a}=b_i^{\mathfrak b}c_i^{\mathfrak b}$
shows that this expression is unchanged. Section values can instead be
read at Gauss radii $r^{1/b_i^{\mathfrak a}}$; Gauss points exist even
when these radii are not in $|K^*|$. We do not identify numerical values
at fixed Gauss points with all gauge-section values.

On the separated normalization each individual branch has radial class
$[-\infty,0]$, so the product class is the whole orthant. Its boundary
map is a bijection of the two branch points, making boundary
fiber-constancy vacuous. But a mixed-zero radial tuple has value
$-\infty$ at one origin and zero at the other; it cannot descend as
an analytic function through their common image. The conductor has
finite value on each branch order, and its zero-dimensional quotient
has no centered normalized section. There is therefore no conflict
with the known conductor quotient theorem on valuation boundaries.
Any operator on all separated radial profiles agreeing with the
degree-one branchwise identity is forced to return that same mixed-zero
vector, outside the connected radial class. This diagnoses a missing
geometric domain condition for P12.6; it does not refute an operator
restricted to compatible source data.

An arbitrary evaluation trace is different. The bounded psh function
\[
 b=\max\{a^{-1}\log|y|+1,-1\}
\]
has value $-1$ at $o$ and values $(-1,0)$ at the same Gauss points,
by GR's cone and maximum properties. On $|x|,|y|\leq r$ it is also
nonpositive. Thus the mixed-zero vector can occur as the evaluation
of a nonradial psh function; our theorem rules out only its literal
homogeneous radial realization.

\section{Comparison, limitations, and evidence}
The catalogue already contains the radial classification, the complete
ideal-orthant envelope, and an unrestricted nodal analytic envelope
counterexample. None is counted as a new resolution. Our additional
calculation is the absence of maximal minorants and a minorant projection
for the specified radial class. P12.3 remains open for an independently
geometric singular envelope theory. P12.2 still lacks the required
general geometric cone/comparison; P12.4 has no products here; P12.5
has no energy here; P12.6 retains its geometric compatibility target.
No exact active atom is solved; the obstruction concerns the stated literal
radial minorant problem.

The exact-rational program \texttt{checks/check\_profiles.py} verifies
340 finite-grid obstacles in dimensions one through four, five node
profiles, and the gauge example. The chosen finite grid has greatest
minorants with zero-face coordinates $-1/4$; its result is explicitly
distinguished from the continuum supremum zero. These are exhaustive
\emph{finite} diagnostics, not proofs of infinite or analytic claims.
The complete plain-text argument is \texttt{P12.3-argument.md}. Source
snapshots, commands, outputs, and environment details accompany the
research artifacts. A separate source-aware Xean audit using
\texttt{gpt-6.1-sol} at \texttt{xhigh} returned \texttt{PASS} for these
precise partial results. Its report and exact primary-source premises are
retained in \texttt{P12.3-review.json}. This is a computational audit, not
a resolution of the original atom or a formal proof certificate.

\begin{thebibliography}{9}
\bibitem{GR} W. Gubler and J. Rabinoff,
\emph{Classically psh and pluriharmonic functions on Berkovich spaces},
arXiv:2506.13548v2, 17 September 2025.
Primary PDF and extracted text preserved in the deliverable.
\bibitem{Book} \emph{Non-Archimedean Pluripotential Theory}, local
book-source repository snapshot, commit
\texttt{00957114db21b627e97369563032ceda9cea09c7}.
Chapter 70, especially Theorems 70.1.5--70.1.10, and Chapter 15's nodal
envelope example. The radial classification is in the included
\texttt{CURVE\_PSH.tex} source; source manifests accompany the draft.
\bibitem{Dossier} Supplied family dossier, \emph{A geometric local theory
on singular germs}, active targets P12.2--P12.6 and their completion
criteria, preserved as \texttt{sources/task-dossier.txt}.
\end{thebibliography}
\end{document}
